\section*{Bab \ref{ch:g10:layered-ocean} --- Samudra yang Berlapis}

\begin{solution}{exo:g10:layered-ocean:1}
\omterm{def:g10:layered-ocean:mixed-layer}{Lapisan tercampur} adalah lapisan permukaan samudra tempat suhu dan salinitas
hampir tidak berubah terhadap kedalaman; \omterm{def:g10:layered-ocean:thermocline}{termoklin} adalah lapisan di bawahnya
tempat suhu turun dengan cepat terhadap kedalaman. Angin dan ombak mengaduk
\omterm{def:g10:layered-ocean:mixed-layer}{lapisan tercampur} (dibantu pada musim dingin oleh pendinginan permukaan).
\end{solution}

\begin{solution}{exo:g10:layered-ocean:2}
\omterm{def:g10:layered-ocean:upwelling}{Penaikan massa air} adalah naiknya air dingin dari bawah lapisan permukaan,
yang menggantikan air permukaan yang didorong menjauh oleh angin. Air yang
naik itu membawa nutrien mineral ke dalam lapisan yang tersinari Matahari,
tempat fitoplankton memakainya untuk tumbuh; fitoplankton menjadi makanan bagi
rantai makanan lainnya.
\end{solution}

\begin{solution}{exo:g10:layered-ocean:3}
Pasifik tropis: sekitar \qty{29}{\degreeCelsius} di permukaan dan
\qty{4.5}{\degreeCelsius} pada \qty{1000}{m}. Samudra Arktik: sekitar
\qty{-1.4}{\degreeCelsius} di permukaan dan mendekati
\qty{0}{\degreeCelsius} pada \qty{1000}{m}.
\end{solution}

\begin{solution}{exo:g10:layered-ocean:4}
Pada \qty{0}{\degreeCelsius} (\qty{1028.1}{kg/m^3} dibandingkan dengan
\qty{1023.3}{kg/m^3} pada \qty{25}{\degreeCelsius}). Air laut: pada
\qty{4}{\degreeCelsius} air murni memiliki massa jenis \qty{1000.0}{kg/m^3},
air laut bersalinitas 35 sekitar \qty{1027.8}{kg/m^3} (dibaca pada grafik).
\end{solution}

\begin{solution}{exo:g10:layered-ocean:5}
Pelampung hanyut pada \qty{1000}{m}, tenggelam sampai \qty{2000}{m}, lalu naik
ke permukaan sambil mengukur suhu, salinitas (dari konduktivitas), dan tekanan
di sepanjang jalan; di permukaan pelampung mengirim profilnya melalui satelit
dan tenggelam lagi. Pelampung hanyut bersama arus ke seluruh samudra dan
bekerja bertahun-tahun, dalam segala musim, jauh dari jalur pelayaran kapal.
\end{solution}

\begin{solution}{exo:g10:layered-ocean:6}
$361.84\times10^6 \times 3.682 \approx \qty{1.33e9}{km^3}$, sekitar
\num{1.33} miliar kilometer kubik. $100/3682 \approx 0.027$: kurang dari
\qty{3}{\percent} kedalaman rata-rata.
\end{solution}

\begin{solution}{exo:g10:layered-ocean:7}
$22.7 - 0.5 = \qty{22.2}{\degreeCelsius}$, dicapai pada \qty{15}{m}: \omterm{def:g10:layered-ocean:mixed-layer}{lapisan
tercampur} sedalam sekitar \qty{15}{m}. Di antara \qty{20}{m} dan \qty{50}{m}:
$(21.8 - 17.3)/30 = \qty{0.15}{\degreeCelsius}$ per meter.
\end{solution}

\begin{solution}{exo:g10:layered-ocean:8}
Matahari hanya memanaskan lapisan permukaan, dan pemanasan serta pendinginan
musiman hanya tersebar di dalam \omterm{def:g10:layered-ocean:mixed-layer}{lapisan tercampur}: di bawah \omterm{def:g10:layered-ocean:thermocline}{termoklin} air
terputus dari keduanya. Pada musim dingin permukaan mendingin, menjadi lebih
rapat, dan tenggelam ke dalam air di bawahnya, dan badai mengaduk laut:
pencampurannya mencapai kedalaman yang besar. Pada musim panas Matahari
menghangatkan lapisan permukaan yang ringan yang mengapung di atas air yang
lebih dingin; angin yang lebih lemah tidak dapat mendorongnya ke bawah, dan
\omterm{def:g10:layered-ocean:mixed-layer}{lapisan tercampur} tetap tipis.
\end{solution}

\begin{solution}{exo:g10:layered-ocean:9}
$0.78/0.18 \approx 4.3$ derajat di dekat \qty{10}{\degreeCelsius};
$0.78/0.06 = 13$ derajat di dekat \qty{0}{\degreeCelsius}. Di laut kutub yang
dingin, perubahan kecil salinitas jauh lebih berpengaruh terhadap massa jenis
daripada perubahan suhu: di sana salinitaslah yang mengendalikan pelapisan.
\end{solution}

\begin{solution}{exo:g10:layered-ocean:10}
Suhunya sama: air yang lebih dalam lebih rapat karena lebih asin (33.3
dibandingkan dengan 28.1). Lapisan di antara keduanya, tempat salinitas
berubah dengan cepat terhadap kedalaman, adalah \omterm{def:g10:layered-ocean:thermocline}{haloklin}; di sini \omterm{def:g10:layered-ocean:thermocline}{haloklin} itu
juga merupakan \omterm{def:g10:layered-ocean:thermocline}{piknoklin}.
\end{solution}

\begin{solution}{exo:g10:layered-ocean:11}
Ketika didorong ke bawah, parsel itu mempertahankan suhu dan salinitasnya,
jadi juga massa jenisnya. Dikelilingi air yang lebih rapat, parsel itu lebih
ringan daripada air yang didesaknya: gaya apung mendorongnya naik kembali ke
arasnya sendiri. Parsel yang didorong ke atas akan lebih rapat daripada
sekitarnya yang baru dan tenggelam kembali. Lapisan-lapisan itu stabil.
\end{solution}

\begin{solution}{exo:g10:layered-ocean:12}
Massa: $0.001 \times 1028.1 = \qty{1.0281}{kg}$. Volume baru: $1.0281/1023.3
\approx \qty{0.0010047}{m^3}$, sekitar \qty{1004.7}{mL}: air itu memuai
sekitar \qty{4.7}{mL}, \qty{0.47}{\percent}. Lapisan setebal \qty{700}{m} akan
menebal sebesar $700 \times 0.0047 \approx \qty{3.3}{m}$.
\end{solution}

\begin{solution}{exo:g10:layered-ocean:13}
Di Atlantik Utara pada bulan Februari. Di sana \qty{150}{m} teratas hampir
sama suhunya, sehingga hampir sama massa jenisnya: mengaduknya hanya
memerlukan sedikit energi. Di Pasifik tropis massa jenis naik sebesar
$1026.0 - 1021.6 =
\qty{4.4}{kg/m^3}$ di antara \qty{50}{m} dan \qty{200}{m}: \omterm{def:g10:layered-ocean:thermocline}{piknoklin} yang
kuat, dan badai harus mendorong air ringan ke bawah ke dalam air yang jauh
lebih rapat.
\end{solution}

\begin{solution}{exo:g10:layered-ocean:14}
Klorofil tinggi jauh dari pesisir: pita di sepanjang khatulistiwa di Pasifik
(dan di Atlantik), tempat angin pasat menggerakkan \omterm{def:g10:layered-ocean:upwelling}{penaikan massa air}; dan
lintang tinggi, seperti Atlantik Utara, tempat \omterm{def:g10:layered-ocean:mixed-layer}{lapisan tercampur} musim dingin
mencapai kedalaman yang besar dan membawa nutrien ke atas. Klorofil sangat
rendah: pusat-pusat pilin subtropis, yang \omterm{def:g10:layered-ocean:stratification}{terstratifikasi} dengan kuat, tempat
tidak ada yang membawa air dalam ke atas dan lapisan yang tersinari Matahari
telah kehabisan nutrien.
\end{solution}

\begin{solution}{exo:g10:layered-ocean:15}
Tenggelam: $1000/(6 \times 3600) \approx \qty{0.046}{m/s}$, sekitar
\qty{4.6}{cm/s}. Naik: $2000/(10 \times 3600) \approx \qty{0.056}{m/s}$,
sekitar \qty{5.6}{cm/s}. $36\,000/2 = 18\,000$ pengukuran, rata-rata berjarak
$2000/18\,000 \approx \qty{0.11}{m}$.
\end{solution}

\begin{solution}{pb:g10:layered-ocean:1}
\textbf{1.} Pilin, sebesar $24.0 - 19.0 = \qty{5.0}{\degreeCelsius}$.

\textbf{2.} Pilin (\qty{20.5}{\degree} LS dibandingkan dengan \qty{13.5}{\degree} LS).
Karena lebih jauh dari khatulistiwa, pilin seharusnya yang lebih sejuk: lintang
saja tidak dapat menjelaskan mengapa pilin \qty{5}{\degreeCelsius} lebih
hangat.

\textbf{3.} $24.0 - 0.5 = \qty{23.5}{\degreeCelsius}$, dicapai di antara
\qty{50}{m} (23.6) dan \qty{100}{m} (22.6), pada sekitar $50 + 50 \times
0.1/1.0 = \qty{55}{m}$.

\textbf{4.} Tidak: suhu turun terus-menerus dari permukaan, sebesar
\qty{1.6}{\degreeCelsius} pada \qty{25}{m} pertama dan
\qty{5.3}{\degreeCelsius} pada \qty{100}{m} pertama; \omterm{def:g10:layered-ocean:thermocline}{termoklin} bermula di
permukaan.

\textbf{5.} Di antara \qty{200}{m} dan \qty{300}{m}: $20.3 - 14.4 =
\qty{5.9}{\degreeCelsius}$ per seratus meter.

\textbf{6.} Di lepas pantai Peru, di antara \qty{50}{m} dan \qty{100}{m}:
sekitar \qty{70}{m} dengan interpolasi (profil lengkapnya memberikan
\qty{65}{m}). Di pilin, di antara \qty{200}{m} dan \qty{300}{m}: sekitar
$200 + 100 \times 5.3/5.9
\approx \qty{290}{m}$.

\textbf{7.} Ya, pada setiap aras kedua kolom: air yang lebih ringan terletak
di atas air yang lebih rapat, dan kedua kolom itu stabil.

\textbf{8.} Air yang lebih asin seharusnya lebih rapat; namun air permukaan
pilin lebih ringan: \qty{1024.71}{kg/m^3} dibandingkan dengan
\qty{1025.08}{kg/m^3}.

\textbf{9.} Air pilin sekitar 1.3 satuan lebih asin, yang menambah sekitar
$1.3 \times 0.8 \approx \qty{1.0}{kg/m^3}$; air itu \qty{5}{\degreeCelsius}
lebih hangat, yang mengurangi sekitar
$5 \times 0.27 \approx \qty{1.4}{kg/m^3}$ pada suhu-suhu ini. Suhulah yang menang.

\textbf{10.} $1026.44 - 1025.08 = \qty{1.36}{kg/m^3}$.

\textbf{11.} $1025.30 - 1024.71 = \qty{0.59}{kg/m^3}$.

\textbf{12.} Dari permukaan sampai \qty{500}{m} massa jenis naik sebesar
$1026.92 - 1024.71 = \qty{2.21}{kg/m^3}$, dan hanya \qty{0.59}{kg/m^3} di
antaranya di atas \qty{200}{m}: sebagian besar, \qty{1.62}{kg/m^3}, terletak
di bawah \qty{200}{m}, di \omterm{def:g10:layered-ocean:thermocline}{termoklin} di antara \qty{200}{m} dan \qty{500}{m}.

\textbf{13.} Di lepas pantai Peru massa jenis naik sebesar $1026.97 - 1025.08 =
\qty{1.89}{kg/m^3}$ dari permukaan sampai \qty{500}{m}, dan
\qty{1.36}{kg/m^3} di antaranya di \qty{200}{m} teratas: sebagian besar
terletak di dekat permukaan.

\textbf{14.} Pada \qty{500}{m} kedua lokasi hampir sama
(\qty{7.9}{\degreeCelsius} dan \qty{7.1}{\degreeCelsius},
\qty{1026.97}{kg/m^3} dan \qty{1026.92}{kg/m^3}): air dingin yang sama
terletak di bawah keduanya; perbedaannya adalah seberapa dekat air itu ke
permukaan.

\textbf{15.} Angin yang bertiup di sepanjang pantai ke arah khatulistiwa
mendorong air permukaan ke laut lepas; air dari bawah naik di dekat pesisir
untuk menggantikannya. Air dingin \omterm{def:g10:layered-ocean:thermocline}{termoklin} tertarik ke atas, sehingga
\omterm{def:g10:layered-ocean:thermocline}{termoklin} miring naik ke arah pesisir dan hampir mencapai permukaan.

\textbf{16.} Di lepas pantai Peru: air bersuhu \qty{13}{\degreeCelsius}
sampai \qty{15}{\degreeCelsius}, dari bawah \omterm{def:g10:layered-ocean:thermocline}{termoklin}, terletak di
\qty{100}{m} teratas, dalam jangkauan pengadukan oleh angin dan jangkauan
cahaya. Di pilin air itu terletak di bawah \qty{250}{m}, di bawah lapisan
hangat dan ringan yang tidak dapat diaduk angin sedalam itu.

\textbf{17.} \omterm{def:g10:layered-ocean:upwelling}{Penaikan massa air} akan berhenti: air permukaan akan menghangat
oleh Matahari, lapisan hangat yang ringan akan terbentuk di atasnya, dan
pasokan nutrien ke lapisan yang tersinari Matahari akan menyusut.

\textbf{18.} Pesisir Peru berwarna biru muda sampai kuning (banyak klorofil);
bagian tengah pilin Pasifik Selatan termasuk daerah berwarna biru paling tua
pada peta itu (sangat sedikit).

\textbf{19.} Fitoplankton memerlukan nutrien selain cahaya. Di pilin,
\omterm{def:g10:layered-ocean:stratification}{stratifikasi} yang kuat menahan air dalam yang kaya nutrien jauh di bawah
lapisan yang tersinari Matahari, dan nutrien di lapisan permukaan telah lama
habis terpakai.

\textbf{20.} $1.36/0.59 \approx 2.3$: \textbf{pertambahan massa jenis pada
\qty{200}{m} teratas sekitar 2.3 kali lebih besar di lepas pantai Peru
daripada di pilin}. \omterm{def:g10:layered-ocean:thermocline}{Piknoklin}, dengan air dingin yang kaya nutrien di
bawahnya, telah terangkat ke dalam \qty{200}{m} teratas di lepas pantai Peru,
sedangkan di pilin \omterm{def:g10:layered-ocean:thermocline}{piknoklin} itu terletak lebih dalam, di bawah \qty{200}{m}.
\end{solution}
